\section{Introduction}
\label{sec:intro}

Weak lensing (WL) is a powerful probe of dark matter and the growth of large-scale structure. In the coming decade, its science cases, including cosmic shear, galaxy-galaxy lensing, and cluster lensing, will reach unprecedented precision with data from Stage-IV surveys. Imaging surveys such as LSST \citep{DESC_SRD}, Roman \citep{Roman_AFTA,HLIS-Multiprobe,HLIS-LSST}, and Euclid \citep{Euclid} will measure the shapes of galaxies over large areas of the sky, while spectroscopic surveys such as DESI \citep{DESI1,DESI2} and SPHEREx \citep{spherex} will provide spectra for tens of millions of galaxies. Much is being done in data analysis and infrastructure to make the most of this vast influx of data. This work represents such an effort as we continue the development of kinematic lensing (KL), a shear measurement technique that reduces shape noise in WL measurements by up to an order of magnitude. Specifically, we explore ways to accelerate KL inference using neural networks.

While weak lensing has been able to achieve good constraints on cosmological parameters \citep{DES_Y3_3x2pt, DESY6CS}, there are numerous sources of statistical uncertainty involved with its measurement. Among these is the so-called shape noise, which is caused by the degeneracy between the intrinsic shape of a galaxy and its shear effect. The intrinsic shape of a galaxy cannot be measured purely by photometry, giving rise to large uncertainties. In WL studies, shape noise is compensated for by averaging over a large number of galaxy samples, but methods to directly reduce shape noise are very limited.

KL was developed with the primary goal of reducing shape noise in individual galaxy measurements \citep{KL}. Using spectroscopic measurements of disk galaxies, KL is able to measure their velocity field, which, when combined with the Tully-Fisher relation (TFR), constrains their intrinsic ellipticity, thereby reducing shape noise. \citep{KLInferenceI} showed with a simulated analysis that KL reduces shape noise from $\sigma_\epsilon^\mathrm{WL}\approx0.26$ to $\sigma_\epsilon^\mathrm{KL}\approx0.022$--$0.041$, and \citep{KLCluster} later made the first detection of a cluster lensing signal with KL, with a typical shear uncertainty of $\sigma(g_+)\approx0.026$, an order of magnitude improvement over WL shape noise. \citep{KLRoman} forecasted that a KL analysis with Roman would improve the figure of merit over a traditional WL analysis by a factor of 1.70 in $\Omega_\mathrm{m}$--$S_8$ and 3.65 in $w_0$--$w_a$. Another forecast for DESI showed that a DESI-KL survey could place highly interesting constraints on structure formation at very low redshift ($z < 0.15$), which would provide unique insights into the $S_8$ tension if it consolidates, while complementing LSST with a different sky footprint \citep{KLDESI}.

However, KL inference currently relies on Markov Chain Monte Carlo (MCMC) sampling of a full forward model, performed separately for each galaxy. In \citep{KLInferenceI}, for example, each fit samples 11 parameters with 100 walkers over 30,000 steps, requiring $3\times10^6$ evaluations of the forward model. This is well suited for analyses of a limited number of galaxies, as was the case for the aforementioned simulated analyses and pilot measurement. However, a DESI-KL survey would contain on the order of $10^7$ galaxies \citep{KLDESI}, each of which would need its own MCMC run. While there have been impressive speed improvements in the pipeline (Everett et al., in prep.), and recent inference runs on DESI datavectors achieved a fitting time of $\sim$20 minutes per galaxy on one GPU (Willett et al., in prep.), performing KL on a full survey would still come at an immense computational cost. Exploration into faster methods of KL inference is therefore needed before KL can be applied at survey scale.

Machine learning methods have been applied to traditional WL shear measurements in order to improve speed. For example, \textsc{Forklens} \citep{Forklens} measures galaxy shapes with a convolutional neural network (CNN) followed by a neural network calibration module. D$_4$CNN$\times$AnaCal \citep{D4CNN} combines an equivariant CNN architecture with analytical shear calibration \citep{Anacal}, and achieves multiplicative biases below percent-level for isolated galaxies. Neural posterior estimation (NPE) has also been used to infer shear directly from simulated multiband images \citep{NPEShear}. To this date, no neural network-based methods have been developed for KL.

In this work, we develop an NPE for KL inference. NPE is a type of simulation-based inference (SBI), which trains neural networks on data generated by a simulator to approximate the posterior (see \citep{SBIGuide} for a review on SBI). Once trained, inference with NPE is amortized: the network can perform inference on new observations without additional training or simulations. Our model uses a CNN to extract features from a galaxy's photometric image and its DESI-like fiber spectra, and a conditional normalizing flow to estimate the joint posterior of shear and galaxy parameters. If successful, this method would greatly accelerate inference for the DESI-KL survey proposed in \citep{KLDESI}. We therefore build our simulations on the DESI and Legacy Survey instrument configurations used in that work, and test our model on simulated populations sampled from the real measured catalog.

We begin with a review of KL theory in Sec.~\ref{sec:kl-basics}. In Sec.~\ref{sec:data-generate} we describe the generation of our training, validation, and test sets. We describe our neural network architecture in Sec.~\ref{sec:nn-arch} and validate the posteriors of our model in Sec.~\ref{sec:results}. In Sec.~\ref{sec:shear}, we explore how the posteriors from our NPE can be turned into shear estimates. We discuss our results in Sec.~\ref{sec:disc}, and conclude in Sec.~\ref{sec:conclusion}.
